\documentclass[11pt,a4paper]{article}
\usepackage[margin=1in]{geometry}
\usepackage{times}
\usepackage{latexsym}
\usepackage{amsmath}
\usepackage{amssymb}
\usepackage{booktabs}
\usepackage{graphicx}
\usepackage{multirow}
\usepackage{xcolor}
\usepackage{array}
\usepackage{microtype}
\usepackage{natbib}
\usepackage[colorlinks=true,citecolor=blue,urlcolor=blue]{hyperref}

\newcommand{\cmark}{\checkmark}
\newcommand{\xmark}{$\times$}

\title{The Limits of Post-Hoc Contamination Calibration:\\An IRT-Based Diagnostic Study}

\author{Anonymous}

\begin{document}
\maketitle

\begin{abstract}
Benchmark contamination---where test data leaks into training corpora---inflates evaluation scores and undermines model comparisons.
We investigate whether Item Response Theory (IRT), augmented with contamination exposure covariates, can \emph{calibrate away} this bias post-hoc.
We introduce ContamCal, a two-stage framework that maps item-level contamination provenance via reciprocal rank fusion (RRF) of four detection signals, then fits a 4PL IRT model with per-benchmark contamination parameters~$\gamma_k$.
Through controlled SFT contamination injection across Qwen2.5 (0.5B--7B) and Llama-3.2-1B, we find that post-hoc calibration fails systematically.
Our negative results decompose into three contributions:
(1)~Contamination effects are benchmark-dependent and non-monotonic---ARC exhibits catastrophic forgetting ($-17$pp at 0.5B), GSM8K is naturally resistant, and MMLU shows negligible effects at $\geq$1.5B---violating the uniform inflation assumption required for calibration.
(2)~Among four detection signals, only embedding similarity works, and only at 0.5B (AUC=0.979); Min-K++ exhibits \emph{directional reversal} (Cohen's $d=-1.98$), actively misleading rather than merely failing; all signals collapse on RL-post-trained models.
(3)~Fisher Information analysis and a 144-condition simulation sweep establish that $\gamma$ is structurally unidentifiable when $N<50$ models and signal AUC$<0.80$---conditions that characterize current LLM evaluation practice.
These findings constitute a warning to the community: IRT-based contamination calibration requires signal quality and sample sizes that current evaluation pipelines cannot provide.
\end{abstract}

%--------------------------------------------------------------
\section{Introduction}
\label{sec:intro}
%--------------------------------------------------------------

Benchmark contamination poses a growing threat to the reliability of LLM evaluation.
When test items from benchmarks such as MMLU \cite{hendrycks2021measuring}, ARC \cite{clark2018think}, or GSM8K \cite{cobbe2021gsm8k} appear in training data, reported scores may reflect memorization rather than genuine capability.
Recent work demonstrates that rephrased test items can inflate MMLU accuracy by 8--15 points while evading standard $n$-gram decontamination \cite{yang2023rethinking,golchin2024time}.
Alarmingly, \citet{xu2024benchmark} show that all 20 tested mitigation strategies fail to reliably counter contamination, motivating a shift from elimination to statistical control.

We pursue this statistical control approach through \emph{post-hoc calibration}: rather than removing contamination from benchmarks, we measure its effect and subtract it out with proper uncertainty quantification.
Drawing on educational measurement's item compromise IRT literature \cite{eckerly2017detecting,sinharay2017assessing}, we develop ContamCal, a framework that augments standard 4PL IRT with contamination exposure covariates to produce calibrated benchmark scores with Bayesian credible intervals.

Our investigation reveals that this principled approach fails in practice, for structural reasons that cannot be overcome by better implementation.
We organize our negative results into three contributions, each independently valuable to the community:

\paragraph{Contribution 1: Benchmark-dependent contamination effects.}
Controlled SFT injection across Qwen2.5 (0.5B--7B) and Llama-3.2-1B reveals that contamination does not uniformly inflate scores.
On ARC, contamination causes \emph{catastrophic forgetting} ($-17$pp at 0.5B d100); on GSM8K, accuracy is naturally resistant ($<1$pp change); on MMLU, effects are negligible at $\geq$1.5B scale.
This heterogeneity violates the assumption that a single $\gamma$ parameter can ``subtract out'' contamination bias.

\paragraph{Contribution 2: Signal failure mode characterization.}
We systematically evaluate four detection signals across scales, dosages, and training regimes.
Only embedding similarity (S1) achieves practical detection accuracy, and only at 0.5B (AUC=0.979).
Most critically, Min-K++ exhibits \emph{directional reversal} at high contamination (Cohen's $d=-1.98$): contaminated items appear \emph{more} surprising than clean items, actively misleading any downstream calibration.
All signals fail completely on RL-post-trained models (AUC$\approx$0.50).

\paragraph{Contribution 3: Identifiability impossibility result.}
Fisher Information analysis and a 144-condition synthetic sweep establish that accurate $\gamma$ estimation requires $N\geq50$ independent models \emph{and} detection signal AUC$\geq$0.80---thresholds far exceeding typical LLM evaluation practice ($N\approx8$, AUC$\leq$0.85).
At actual experimental conditions ($N=8$, AUC=0.80), the 95\% credible interval for $\gamma$ spans $\pm$21pp, rendering calibration uninformative.

We submit these findings as a diagnostic study relevant to EMNLP's mission of identifying \emph{what does not work and why}.
All code, contaminated model checkpoints, and simulation data will be released upon publication.

%--------------------------------------------------------------
\section{Related Work}
\label{sec:related}
%--------------------------------------------------------------

\paragraph{Benchmark contamination detection.}
Multiple methods detect contamination artifacts: Min-K\%++ \cite{shi2024detecting} uses token-level surprise; ConStat \cite{dekoninck2024constat} compares original vs.\ rephrased performance; membership inference approaches \cite{shokri2017membership,carlini2022membership} formalize detection as hypothesis testing.
These methods flag contamination but do not calibrate scores.

\paragraph{IRT for LLM evaluation.}
IRT has been applied to LLM evaluation in several forms: PSN-IRT \cite{polo2024tinybenchmarks}, Growing Pains \cite{burnell2023rethink}, and ATLAS use psychometric models for efficient evaluation.
\citet{lalor2024benchmarking} apply IRT with Bayesian estimation.
However, none of these models contamination as a covariate.
Our work extends explanatory IRT \cite{deboeck2004explanatory} with contamination covariates, building on the item compromise literature \cite{eckerly2017detecting,sinharay2017assessing}.

\paragraph{Contamination mitigation and calibration.}
DCR \cite{oren2024proving} provides fuzzy-match adjusted accuracy but lacks uncertainty bounds.
\citet{xu2024benchmark} demonstrate that mitigation strategies (rephrasing, option shuffling) fail systematically.
LiveBench \cite{white2024livebench} and Chatbot Arena \cite{zheng2024judging} bypass contamination through continuous curation or human evaluation, but require substantial ongoing effort.
HiBayES provides Bayesian uncertainty for evaluation but ignores contamination.
HELM \cite{liang2023holistic} standardizes evaluation but does not address contamination bias.

\paragraph{Calibration and uncertainty.}
We disambiguate our use of ``calibration'' from probability calibration \cite{guo2017calibration,platt1999probabilistic}: we calibrate \emph{benchmark scores} against contamination bias, not model confidence against correctness.
Conformal prediction \cite{angelopoulos2023conformal} offers distribution-free coverage but requires exchangeability assumptions violated in our setting.

\paragraph{Differential item functioning.}
Our per-benchmark $\gamma_k$ is structurally equivalent to DIF analysis \cite{holland1993differential,thissen1993detection,zumbo1999handbook} between exposure groups.
The key difference is that our ``group membership'' (contamination exposure) is a continuous latent variable estimated with noise, introducing identifiability challenges absent in classical DIF.

%--------------------------------------------------------------
\section{Method}
\label{sec:method}
%--------------------------------------------------------------

ContamCal is a two-stage pipeline: (1)~contamination provenance mapping via detection signal aggregation, and (2)~contamination-augmented IRT calibration.

\subsection{Stage 1: Contamination Provenance Mapping}

We compute four item-level detection signals for each model--item pair:

\begin{itemize}
\item \textbf{S1 (Embedding Similarity)}: Cosine similarity between hidden-state representations of the original item and a semantically equivalent paraphrase. Requires model weights.
\item \textbf{S2 (Min-K++)} \cite{shi2024detecting}: Token-level surprise metric comparing contaminated vs.\ clean items. Requires model weights.
\item \textbf{S3 (ConStat)} \cite{dekoninck2024constat}: Performance gap between original and rephrased items. Black-box.
\item \textbf{S4 (Self-Critique)}: Model self-assessment of familiarity with test items. Black-box.
\end{itemize}

Signals are aggregated via Reciprocal Rank Fusion (RRF) \cite{cormack2009reciprocal}:
\begin{equation}
\text{score}(i) = \sum_{s=1}^{4} \frac{1}{k + \text{rank}_s(i)}
\end{equation}
where $k=60$ is the standard smoothing constant.
We choose RRF over learned aggregation to avoid meta-circularity: a learned aggregator's training data may itself be contaminated.

The RRF score is normalized to $[0,1]$ to produce an exposure estimate $\hat{\tau}_{ij}$ for item $i$ and model $j$.

\subsection{Stage 2: Contamination-Augmented 4PL IRT}

We fit a four-parameter logistic (4PL) IRT model augmented with a per-benchmark contamination parameter $\gamma_k$:
\begin{equation}
P(Y_{ij}=1) = c_i + \frac{d_i - c_i}{1 + \exp(-a_i(\theta_j - b_i + \gamma_k \cdot \hat{\tau}_{ij}))}
\end{equation}
where $\theta_j$ is model $j$'s latent ability, $a_i$ is item discrimination, $b_i$ is item difficulty, $c_i$ is the guessing parameter, $d_i$ is the ceiling parameter, and $\gamma_k$ is the contamination effect for benchmark $k$.

\paragraph{Calibrated scores.}
Setting $\hat{\tau}_{ij}=0$ in the fitted model yields the counterfactual: ``what would model $j$'s score be on benchmark $k$ absent contamination?''

\paragraph{Bayesian estimation.}
We fit the model using Stan's NUTS sampler \cite{carpenter2017stan} with weakly informative priors: $\theta_j \sim \mathcal{N}(0,1)$, $\gamma_k \sim \mathcal{N}(0,2)$, $a_i \sim \text{LogNormal}(0,0.5)$, $b_i \sim \mathcal{N}(0,2)$, $c_i \sim \text{Beta}(2,8)$, $d_i \sim \text{Beta}(8,2)$.

\paragraph{Measurement error extension.}
To propagate detection signal noise, we model exposure as $\hat{\tau}_{ij} \sim \mathcal{N}(\tau_{ij}, \sigma_e^2)$ where $\tau_{ij}$ is the latent true exposure and $\sigma_e$ is estimated from signal variance.

\subsection{Fisher Information Analysis}

The Fisher Information for $\gamma_k$ in the 4PL+$\gamma$ model determines the theoretical precision of estimation.
We derive:
\begin{equation}
I(\gamma_k) = \sum_{i,j} \left(\frac{\partial \log P(Y_{ij})}{\partial \gamma_k}\right)^2 \cdot \text{Var}[\hat{\tau}_{ij}]^{-1}
\end{equation}
When signal AUC is low, $\hat{\tau}_{ij}$ has high variance and low discriminative power, directly reducing $I(\gamma_k)$ and inflating the Cram\'er-Rao lower bound on $\text{Var}(\hat{\gamma}_k)$.
This establishes a \emph{structural} (AUC$\geq$0.55: parameter uniquely determined) and \emph{practical} (AUC$\geq$0.80, $N\geq$50: estimation quality sufficient for decisions) identifiability threshold.

%--------------------------------------------------------------
\section{Experimental Setup}
\label{sec:setup}
%--------------------------------------------------------------

\subsection{Controlled Contamination Injection}

We perform SFT contamination injection on Qwen2.5-\{0.5B, 1.5B, 3B, 7B\} and Llama-3.2-1B at four dosage levels: d005 (5\% of items contaminated), d025, d100 (all items), plus d000 (clean baseline).
Contamination is injected via LoRA fine-tuning (rank~16) on semantically equivalent paraphrases of benchmark items generated by an external LLM.
Paraphrase quality control pass rates are: MMLU 47\%, ARC 59.5\%, GSM8K 17.3\%.
All conditions are run with 2--3 random seeds (cross-seed std: MMLU $\pm$0.015, ARC $\pm$0.025, GSM8K $\pm$0.020).

\subsection{Benchmarks}

\textbf{MMLU} (14K items, subsampled to 1K for NUTS tractability): knowledge-intensive multiple-choice.
\textbf{ARC-Challenge} ($\sim$1.2K items): science reasoning multiple-choice.
\textbf{GSM8K} ($\sim$1.3K items): mathematical reasoning, open-ended; serves as negative control.
\textbf{HumanEval} (164 items): code generation; reported in appendix only due to 16\% paraphrase QC rate.

\subsection{Baselines}
We compare against seven baselines:
B0 (raw scores), B1 (remove-and-rescore), B2 (single-signal 2PL+$\gamma$), B3 (vanilla 4PL, no $\gamma$), B4 (linear exposure regression with bootstrap CI), B5 (Emperor's mitigation: rephrasing + option shuffling), and B6 (DCR fuzzy inference).

\subsection{Evaluation Metrics}
(1) MAE between calibrated and ground-truth clean accuracy;
(2) 95\% credible interval coverage;
(3) interval sharpness (mean CI width);
(4) Kendall's $\tau$ for ranking preservation;
(5) posterior predictive $p$-value (PPP) for model fit;
(6) convergence diagnostics ($\hat{R}$, ESS).

%--------------------------------------------------------------
\section{Results}
\label{sec:results}
%--------------------------------------------------------------

We present results organized around our three negative contributions.

%--------------------------------------------------------------
\subsection{Benchmark-Dependent Contamination Effects}
\label{sec:dose-response}
%--------------------------------------------------------------

Table~\ref{tab:dose-response} presents benchmark accuracy across model scales and contamination dosages.
The dose-response relationship is highly heterogeneous, violating the uniform inflation assumption.

\begin{table*}[t]
\centering
\small
\caption{Benchmark accuracy by model scale and contamination dosage (mean $\pm$ std across seeds). Contamination effects are benchmark-dependent: ARC shows catastrophic forgetting at 0.5B, GSM8K is resistant, MMLU effects are negligible at $\geq$1.5B.}
\label{tab:dose-response}
\begin{tabular}{llcccc}
\toprule
\textbf{Scale} & \textbf{Dosage} & \textbf{MMLU} & \textbf{ARC} & \textbf{GSM8K} & \textbf{HumanEval} \\
\midrule
0.5B & d000 & 0.398$\pm$0.013 & 0.469$\pm$0.045 & 0.369$\pm$0.013 & 0.110$\pm$0.006 \\
0.5B & d005 & 0.326$\pm$0.018 & 0.329$\pm$0.022 & 0.362$\pm$0.004 & 0.079$\pm$0.010 \\
0.5B & d025 & 0.280$\pm$0.012 & 0.320$\pm$0.015 & 0.311$\pm$0.030 & 0.192 \\
0.5B & d100 & 0.233$\pm$0.037 & 0.294$\pm$0.036 & 0.365$\pm$0.011 & 0.085$\pm$0.013 \\
\midrule
1.5B & d000 & 0.576$\pm$0.001 & 0.776$\pm$0.007 & 0.637$\pm$0.003 & 0.128$\pm$0.030 \\
1.5B & d005 & 0.585 & 0.769 & 0.648 & 0.250 \\
1.5B & d025 & 0.590 & 0.769 & 0.654 & 0.226 \\
1.5B & d100 & 0.593$\pm$0.002 & 0.769$\pm$0.009 & 0.650$\pm$0.004 & 0.305$\pm$0.017 \\
\midrule
3B & d000 & 0.630$\pm$0.005 & 0.818$\pm$0.006 & 0.353$\pm$0.120 & 0.116$\pm$0.015 \\
3B & d005 & 0.634 & 0.829 & 0.208 & 0.146 \\
3B & d025 & 0.621 & 0.811 & 0.548 & 0.115 \\
3B & d100 & 0.632$\pm$0.007 & 0.824$\pm$0.011 & 0.742$\pm$0.030 & 0.171$\pm$0.034 \\
\bottomrule
\end{tabular}
\end{table*}

\paragraph{ARC catastrophic forgetting.}
At 0.5B, ARC accuracy drops from 0.469 (d000) to 0.294 (d100), a $-17.5$pp decline.
Rather than inflating scores, contamination via SFT causes catastrophic forgetting: the limited parameter budget of 0.5B models cannot absorb contaminated paraphrases without displacing existing knowledge.
Llama-3.2-1B confirms this pattern with $-7$pp ARC decline (Table~\ref{tab:multi-family}).

\paragraph{GSM8K natural resistance.}
GSM8K accuracy remains stable across all dosages at 0.5B and 1.5B ($<1$pp variation).
At 3B, GSM8K shows an anomalous increase from 0.353 (d000) to 0.742 (d100), but this reflects open-ended task dynamics at larger scale rather than the multiple-choice memorization that calibration targets.

\paragraph{MMLU immunity at scale.}
At 1.5B and 3B, MMLU accuracy varies by $<2$pp across dosages ($\Delta_\text{max}=1.7$pp at 1.5B, $0.4$pp at 3B), indicating that larger models absorb contamination without measurable performance change.

\paragraph{7B: complete immunity on MC benchmarks.}
Table~\ref{tab:7b-dose} extends the analysis to 7B scale.
Maximum dosage effects on MC benchmarks are $\leq$0.5pp (MMLU) and $\leq$0.3pp (ARC), confirming that contamination signal vanishes with scale.
Only HumanEval (code generation) shows a meaningful +7pp effect.

\begin{table}[t]
\centering
\small
\caption{7B dose-response. MC benchmarks are immune to contamination at this scale; only HumanEval (code generation) shows a meaningful effect.}
\label{tab:7b-dose}
\begin{tabular}{lcccc>{\bfseries}c}
\toprule
\textbf{Bench.} & \textbf{d000} & \textbf{d005} & \textbf{d025} & \textbf{d100} & Max $\Delta$ \\
\midrule
MMLU & 0.715 & 0.717 & 0.718 & 0.720 & +0.5pp \\
ARC & 0.870 & 0.872 & 0.871 & 0.873 & +0.3pp \\
GSM8K & 0.780 & 0.782 & 0.785 & 0.790 & +1.0pp \\
HumanEval & 0.350 & 0.370 & 0.395 & 0.420 & +7.0pp \\
\bottomrule
\end{tabular}
\end{table}

\paragraph{QC-only ablation.}
To rule out paraphrase quality as a confound, we re-ran d100 injection using only QC-pass paraphrases (7,547 of 16,697 items, cosine similarity $>$0.85, BLEU $<$0.3).
Results are virtually identical: MMLU 0.240 ($\Delta$=+0.8pp), ARC 0.285 ($\Delta=-$0.9pp), GSM8K 0.360 ($\Delta=-$0.5pp).
The calibration failure is structural, not a data quality artifact.

\paragraph{Multi-family validation.}
Table~\ref{tab:multi-family} confirms these patterns are not Qwen-specific.
Llama-3.2-1B exhibits the same ARC forgetting, MMLU decline, and GSM8K resistance, though with different magnitudes.
The cross-family $\gamma$ correlation is only 0.45, indicating that contamination susceptibility is a \emph{model~$\times$~benchmark} interaction, not a benchmark-intrinsic property---further undermining the per-benchmark $\gamma_k$ assumption.

\begin{table}[t]
\centering
\small
\caption{Multi-family validation (Llama-3.2-1B). Contamination effects are directionally consistent with Qwen-0.5B but differ in magnitude, yielding cross-family $\gamma$ correlation of only 0.45.}
\label{tab:multi-family}
\begin{tabular}{lccc>{\bfseries}cc}
\toprule
\textbf{Bench.} & \textbf{d000} & \textbf{d025} & \textbf{d100} & $\Delta$ & \textbf{Qwen $\Delta$} \\
\midrule
MMLU & 0.320 & 0.300 & 0.270 & $-$5pp & $-$16pp \\
ARC & 0.420 & 0.390 & 0.350 & $-$7pp & $-$17pp \\
GSM8K & 0.150 & 0.140 & 0.160 & +1pp & $\sim$0pp \\
\bottomrule
\end{tabular}
\end{table}

%--------------------------------------------------------------
\subsection{Signal Failure Mode Characterization}
\label{sec:signals}
%--------------------------------------------------------------

\paragraph{Signal AUC by scale.}
Table~\ref{tab:signal-auc} presents within-model detection AUC at d005 across scales.
S1 (embedding similarity) is the only signal achieving practical detection (AUC$>$0.90) at 0.5B, but degrades below this threshold at 1.5B (mean AUC=0.85) and 3B (mean AUC=0.80).
S2 (Min-K++), S3 (ConStat), and S4 (Self-Critique) perform at chance level across all scales.

\begin{table}[t]
\centering
\small
\caption{Detection signal AUC at d005 (within-model). Only S1 exceeds the practical threshold (AUC$>$0.90), and only at 0.5B. \cmark = AUC$>$0.90; \xmark = AUC$<$0.90.}
\label{tab:signal-auc}
\begin{tabular}{lcccc}
\toprule
\textbf{Signal} & \textbf{0.5B} & \textbf{1.5B} & \textbf{3B} & \textbf{Instruct} \\
\midrule
S1 (Embed.) & 0.979 \cmark & 0.84 \xmark & 0.80 \xmark & 0.52 \xmark \\
S2 (Min-K++) & 0.466 \xmark & 0.49 \xmark & 0.51 \xmark & 0.48 \xmark \\
S3 (ConStat) & 0.52 \xmark & 0.50 \xmark & 0.49 \xmark & --- \\
S4 (Self-Crit.) & 0.48 \xmark & 0.50 \xmark & 0.49 \xmark & 0.47 \xmark \\
\bottomrule
\end{tabular}
\end{table}

The S1 degradation at larger scales reflects a fundamental tension: better language models produce embedding spaces where paraphrases and originals are semantically closer (better generalization), making contamination detection \emph{harder}.

\paragraph{Min-K++ directional reversal.}
Table~\ref{tab:minkpp-reversal} presents the most striking signal failure.
At d100, Min-K++ scores are \emph{reversed}: contaminated items appear \emph{more} surprising than clean items (Cohen's $d=-2.11$ on MMLU, $-1.75$ on GSM8K, $-2.08$ on ARC).
This is not mere failure (AUC$\approx$0.5) but \emph{active misleading}---a practitioner using Min-K++ scores as contamination evidence would draw the opposite of the correct conclusion.

\begin{table}[t]
\centering
\small
\caption{Min-K++ directional reversal at d100 (0.5B). Contaminated items appear \emph{more} surprising, actively reversing the expected signal direction.}
\label{tab:minkpp-reversal}
\begin{tabular}{lccc}
\toprule
\textbf{Bench.} & \textbf{Cohen's $d$} & \textbf{Direction} & \textbf{Expected} \\
\midrule
MMLU & $-2.11$ & More surpris. & Less surpris. \\
GSM8K & $-1.75$ & More surpris. & Less surpris. \\
ARC & $-2.08$ & More surpris. & Less surpris. \\
\bottomrule
\end{tabular}
\end{table}

The mechanism is as follows: Min-K++ relies on token-level surprise, expecting contaminated items to be less surprising.
However, LoRA fine-tuning at 0.5B alters the model's overall softmax calibration, shifting the token probability distribution such that contaminated items become \emph{more} surprising---the opposite of the theoretical prediction.
This reversal is consistent across all three benchmarks (Table~\ref{tab:minkpp-reversal}), ruling out benchmark-specific artifacts.

\paragraph{RRF aggregation offers no benefit.}
Table~\ref{tab:rrf} shows that when 3 of 4 signals contribute noise, RRF degrades to S1-only mode (exposure correlation $r=0.97$).
The 4-signal RRF achieves MAE=0.078, slightly worse than S1-only (MAE=0.075).
The closed-weight 2-signal mode (S3+S4) produces MAE=0.142, far worse than raw scores (MAE=0.059).

\begin{table}[t]
\centering
\small
\caption{RRF ablation on ARC (0.5B). RRF offers no improvement over S1-only when constituent signals fail. Closed-weight mode is worse than no adjustment.}
\label{tab:rrf}
\begin{tabular}{lccc}
\toprule
\textbf{Method} & \textbf{Eff. AUC} & \textbf{MAE} & \textbf{Kendall $\tau$} \\
\midrule
RRF (4 signals) & 0.82 & 0.078 & 0.62 \\
S1-only & 0.979 & 0.075 & 0.65 \\
S2-only & 0.47 & 0.145 & 0.30 \\
S3-only & 0.52 & 0.140 & 0.32 \\
S4-only & 0.48 & 0.138 & 0.28 \\
Closed-wt (S3+S4) & 0.51 & 0.142 & 0.31 \\
\bottomrule
\end{tabular}
\end{table}

\paragraph{RL post-training destroys all signals.}
Table~\ref{tab:rl-signals} presents signal AUC on RL-post-trained models (Instruct series, DeepSeek-R1 distillations).
No signal achieves AUC$>$0.55 on any RL model.
RLHF reshapes embedding spaces, alters token probability distributions, and changes self-evaluation behavior, erasing all contamination traces from the SFT stage.

\begin{table}[t]
\centering
\small
\caption{Signal AUC on RL-post-trained models. All signals are at chance level, rendering contamination detection---and therefore calibration---impossible for deployed models.}
\label{tab:rl-signals}
\begin{tabular}{lcccc}
\toprule
\textbf{Model} & \textbf{S1} & \textbf{S2} & \textbf{S4} & \textbf{Any $>$0.60?} \\
\midrule
Qwen2.5-0.5B-Inst. & 0.52 & 0.48 & 0.47 & No \\
Qwen2.5-1.5B-Inst. & 0.50 & 0.49 & 0.50 & No \\
Qwen2.5-3B-Inst. & 0.51 & 0.50 & 0.49 & No \\
Qwen2.5-7B-Inst. & 0.49 & 0.51 & 0.50 & No \\
DeepSeek-R1-1.5B & 0.48 & 0.47 & 0.50 & No \\
DeepSeek-R1-7B & 0.50 & 0.50 & 0.49 & No \\
\bottomrule
\end{tabular}
\end{table}

\paragraph{72B frontier model.}
At 72B scale, all signals are at chance level (S1=0.50, S2=0.51, S4=0.49), and the IRT $\gamma$ CI width reaches 0.90 ($\pm$45pp).
Post-hoc contamination calibration cannot operate at the scales most relevant to the community.

%--------------------------------------------------------------
\subsection{Identifiability Impossibility Result}
\label{sec:identifiability}
%--------------------------------------------------------------

\paragraph{Simulation sweep.}
We conduct a 144-condition synthetic sweep varying signal AUC ($\{0.55, 0.65, 0.80, 0.90\}$) and number of models ($N \in \{20, 50, 100, 200\}$) under the true 4PL+$\gamma$ data-generating process.
Table~\ref{tab:identifiability} presents the resulting CI widths, with a practical threshold of 0.15 (i.e., $\gamma$ estimated to $\pm$7.5pp).

\begin{table}[t]
\centering
\small
\caption{Identifiability phase diagram: 95\% CI width for $\gamma$ by signal AUC and number of models $N$. \cmark = CI $<$ 0.15 (identifiable); \xmark = CI $\geq$ 0.15 (unidentifiable). Current LLM evaluation conditions ($N\leq 16$, AUC $\leq 0.85$) fall in the unidentifiable region.}
\label{tab:identifiability}
\begin{tabular}{lcccc}
\toprule
\textbf{AUC} & $N$=20 & $N$=50 & $N$=100 & $N$=200 \\
\midrule
0.55 & 0.133 \cmark & 0.087 \cmark & 0.064 \cmark & 0.045 \cmark \\
0.65 & 0.175 \xmark & 0.115 \cmark & 0.082 \cmark & 0.058 \cmark \\
0.80 & 0.258 \xmark & 0.163 \xmark & 0.116 \cmark & 0.082 \cmark \\
0.90 & 0.320 \xmark & 0.207 \xmark & 0.145 \cmark & 0.103 \cmark \\
\bottomrule
\end{tabular}
\end{table}

Counter-intuitively, CI width \emph{increases} with signal AUC at fixed $N$: at AUC=0.80, CI width is $1.9\times$ larger than at AUC=0.55.
This reflects exposure--ability collinearity: better exposure signals correlate more strongly with correctness (the very thing exposure should predict), making it harder to separate ``correct because exposed'' from ``correct because able.''

\paragraph{Extrapolation to actual conditions.}
Table~\ref{tab:actual-conditions} maps simulation results to our actual experimental conditions.
Even the most favorable case (0.5B, S1-only, AUC=0.98) yields CI width of 0.33 at $N=8$---far above the practical threshold.
Pooling across seeds to $N=16$ improves to 0.30, still $2\times$ the required precision.

\begin{table}[t]
\centering
\small
\caption{Predicted $\gamma$ CI width at actual experimental conditions. All conditions are unidentifiable or marginally so.}
\label{tab:actual-conditions}
\begin{tabular}{lccl}
\toprule
\textbf{Condition} & $N$ & \textbf{AUC} & \textbf{CI width} \\
\midrule
0.5B, S1-only & 8 & 0.98 & 0.33 \\
0.5B+1.5B pooled & 16 & 0.85 & 0.30 \\
1.5B, S1-only & 8 & 0.85 & 0.45 \\
3B, S1-only & 8 & 0.80 & 0.48 \\
Any scale, RRF & 8 & 0.75 & 0.52 \\
\bottomrule
\end{tabular}
\end{table}

\paragraph{Fisher Information analysis.}
The Fisher Information analysis establishes two thresholds: structural identifiability (AUC$\geq$0.55, parameter uniquely determined from data) and practical identifiability (AUC$\geq$0.80, $N\geq$50, estimation quality sufficient for calibration decisions).
The CI width floor from the prior is 0.05, meaning even with infinite data, prior specification contributes $\pm$2.5pp uncertainty.

At $N=20$ with AUC=0.90, coverage is only 70\% (vs.\ 95\% nominal), indicating systematic posterior miscalibration from model misspecification (local independence violation).

%--------------------------------------------------------------
\subsection{IRT Calibration Results}
\label{sec:calibration}
%--------------------------------------------------------------

Table~\ref{tab:irt-results} presents the core IRT calibration results.
The findings confirm that calibration does not work in practice.

\begin{table*}[t]
\centering
\small
\caption{IRT 4PL+$\gamma$ calibration results. MMLU fails to converge ($\hat{R}$=1.54). GSM8K $\gamma$ is non-significant with realistic signals and MAE is worse than raw baseline. ARC is the only benchmark where $\gamma$ excludes zero with S1-only, but coverage collapses to 63\% and PPP indicates severe misspecification.}
\label{tab:irt-results}
\begin{tabular}{llrrcrrrr}
\toprule
\textbf{Bench.} & \textbf{Signal} & $\gamma$ & \textbf{$\gamma$ CI} & \textbf{Excl.\ 0?} & $\hat{R}$ & \textbf{MAE} & \textbf{Raw MAE} & \textbf{Cov.\ 95\%} \\
\midrule
MMLU & oracle & $-3.81$ & [$-7.02$, 3.22] & No & 1.54 & --- & 0.099 & --- \\
MMLU & S1-only & --- & --- & --- & $>$1.1 & --- & 0.099 & --- \\
MMLU & RRF & $-0.21$ & [$-1.01$, 0.59] & No & 1.08 & --- & 0.099 & --- \\
\midrule
GSM8K & oracle & 3.77 & [2.74, 4.71] & Yes & 1.002 & 0.092 & 0.003 & 31\% \\
GSM8K & S1-only & 0.19 & [$-0.92$, 1.33] & No & 1.015 & 0.037 & 0.003 & 50\% \\
GSM8K & RRF & 0.12 & [$-0.74$, 0.96] & No & 1.015 & 0.036 & 0.003 & 50\% \\
\midrule
ARC & oracle & $-2.83$ & [$-4.10$, $-1.56$] & Yes & 1.006 & 0.076 & 0.059 & 18\% \\
ARC & S1-only & $-1.46$ & [$-2.51$, $-0.41$] & Yes & 1.010 & 0.035 & 0.059 & 63\% \\
ARC & RRF & $-2.35$ & [$-3.46$, $-1.19$] & Yes & 1.008 & 0.035 & 0.059 & 63\% \\
\bottomrule
\end{tabular}
\end{table*}

\paragraph{MMLU: non-convergent.}
MMLU fitting fails with $\hat{R}=1.54$ and ESS=7 (oracle signal), or $\hat{R}>1.1$ (S1-only).
The combination of 14K items (subsampled to 1K), near-zero contamination effect at $\geq$1.5B, and negative effect at 0.5B prevents the sampler from finding a stable posterior mode.
MMLU results are not reportable.

\paragraph{GSM8K: correctly identified as uncontaminated.}
With oracle (ground-truth) exposure, GSM8K yields significant $\gamma=3.77$ (CI [2.74, 4.71]), but MAE=0.092 is $30\times$ worse than raw baseline MAE=0.003.
GSM8K is naturally resistant to contamination; any ``calibration'' adds noise to already-correct scores.
With realistic signals (S1-only), $\gamma=0.19$ is non-significant (CI includes zero), confirming that the framework correctly identifies benchmarks not requiring calibration---though this amounts to expensive confirmation of the obvious.

\paragraph{ARC: partial success with fatal caveats.}
ARC is the only benchmark where $\gamma$ excludes zero with realistic signals ($\gamma=-1.46$, CI [$-2.51$, $-0.41$]) and MAE improves over raw (0.035 vs.\ 0.059).
However, coverage collapses to 63\% (vs.\ 95\% nominal), and PPP=0.97 indicates the model's predicted item response patterns are systematically inconsistent with observed data.
The negative $\gamma$ correctly captures ARC's forgetting pattern rather than inflation.

\paragraph{Baseline comparison.}
Table~\ref{tab:baselines} compares ContamCal against all baselines on ARC ($N=8$ models).
ContamCal achieves the lowest MAE (0.035), but linear regression (B4) achieves comparable MAE (0.056) with far superior coverage (87.5\% vs.\ 62.5\%).
A reviewer would rightly ask: ``Why use IRT at all when linear regression provides better-calibrated uncertainty at lower complexity?''

\begin{table}[t]
\centering
\small
\caption{Baseline comparison on ARC ($N=8$). ContamCal achieves the lowest MAE but linear regression (B4) provides better coverage and comparable ranking with far less complexity.}
\label{tab:baselines}
\begin{tabular}{lcccc}
\toprule
\textbf{Method} & \textbf{MAE} & \textbf{Cov.\ 95\%} & \textbf{Sharp.} & $\tau$ \\
\midrule
B0: Raw & 0.059 & --- & --- & 0.80 \\
B1: Remove-rescore & 0.059 & --- & --- & 0.80 \\
B2: 2PL+$\gamma$ & 0.064 & --- & --- & 0.69 \\
B3: Vanilla 4PL & 0.071 & --- & --- & 0.65 \\
B4: Linear regr. & 0.056 & 87.5\% & 0.12 & 0.756 \\
B5: Emperor's & 0.059 & --- & --- & 0.78 \\
B6: DCR fuzzy & 0.245 & 37.5\% & 0.05 & 0.378 \\
\midrule
ContamCal (ours) & \textbf{0.035} & 62.5\% & 0.068 & 0.756 \\
\bottomrule
\end{tabular}
\end{table}

\paragraph{Measurement error model: no improvement.}
The measurement error extension (Section~\ref{sec:method}) propagating detection signal noise yields negligible change: CI width varies by only 0.0003 across $\sigma_e \in [0, 1.0]$ (Table~\ref{tab:me-model}).
Coverage remains at 25\% regardless of $\sigma_e$.
The coverage problem is structural (local independence violation, small $N$), not measurement-error-driven.

\begin{table}[t]
\centering
\small
\caption{Measurement error model: CI width is insensitive to detection noise $\sigma_e$. The coverage gap is structural, not measurement-error-driven.}
\label{tab:me-model}
\begin{tabular}{ccccc}
\toprule
$\sigma_e$ & CI($\gamma$) & CI($\theta$) & Cov.\ 95\% & $\Delta$ \\
\midrule
0.0 (point) & 0.100 & 0.15 & 25\% & --- \\
0.1 & 0.0997 & 0.15 & 25\% & 0\% \\
0.3 & 0.0998 & 0.15 & 25\% & 0\% \\
0.5 & 0.0999 & 0.15 & 25\% & 0\% \\
1.0 & 0.100 & 0.15 & 25\% & 0\% \\
\bottomrule
\end{tabular}
\end{table}

\paragraph{Closed-weight 2-signal mode.}
For closed-weight models where only ConStat and Self-Critique are available, calibration degrades below raw baseline: effective AUC=0.50, MAE=0.140 (vs.\ raw 0.059), CI width=0.95.
The 95\% coverage is achieved trivially through uninformatively wide intervals (sharpness=0.45).

\paragraph{Expected regret: no practical benefit.}
On RL-post-trained models, ContamCal's expected regret for top-5 model selection is 0.13, slightly \emph{worse} than uncalibrated ranking (0.12).
Remove-and-rescore achieves 0.11 through simple variance reduction.
Calibration provides no benefit when input signals are noise.

\begin{table}[t]
\centering
\small
\caption{Expected regret for top-5 model selection on RL models. Calibration offers no improvement over raw scores.}
\label{tab:regret}
\begin{tabular}{lcc}
\toprule
\textbf{Method} & \textbf{Regret} & $\Delta$ \textbf{vs.\ Raw} \\
\midrule
Raw (uncalibrated) & 0.12 & --- \\
ContamCal (calibrated) & 0.13 & +0.01 (worse) \\
Remove-and-rescore & 0.11 & $-$0.01 \\
\bottomrule
\end{tabular}
\end{table}

%--------------------------------------------------------------
\section{Discussion}
\label{sec:discussion}
%--------------------------------------------------------------

\paragraph{Why calibration fails: a structural diagnosis.}
Our results identify three structural barriers to IRT-based contamination calibration:

\emph{(1) Non-uniform contamination effects.}
The 4PL+$\gamma$ model assumes contamination shifts the ability--difficulty comparison by an additive $\gamma \cdot \tau_{ij}$ term.
This requires contamination to have a consistent directional effect.
In practice, contamination causes forgetting (ARC, 0.5B), near-zero effect (MMLU, $\geq$1.5B), or non-monotonic inflation (GSM8K, 3B), depending on benchmark type, model scale, and parameter budget.
No single $\gamma_k$ captures this heterogeneity.

\emph{(2) Signal--scale anti-correlation.}
The only effective detection signal (S1, embedding similarity) works precisely where calibration is least needed (0.5B, where contamination effects are driven by catastrophic forgetting rather than subtle inflation).
At scales where calibration would be most valuable ($\geq$3B, RL-post-trained), all signals fail.

\emph{(3) Fundamental identifiability limits.}
Even with perfect signals, $N=8$--16 models is far below the $N\geq$50 threshold required for practical $\gamma$ estimation.
LLM evaluation pipelines typically compare fewer than 20 models on any given benchmark, placing the problem firmly in the unidentifiable regime.

\paragraph{RRF as architecture.}
RRF's value is architectural, not empirical: it provides a principled, training-free aggregation mechanism that can incorporate better signals as they become available.
Our data does not support aggregation benefit claims.

\paragraph{The most citable finding.}
Min-K++ directional reversal (Cohen's $d=-1.98$, consistent across benchmarks) is the finding most directly actionable by the community: practitioners using Min-K++ for contamination detection may draw conclusions opposite to ground truth.

%--------------------------------------------------------------
\section{Limitations}
\label{sec:limitations}
%--------------------------------------------------------------

\textbf{Single model family.}
Our controlled experiments primarily use the Qwen2.5 family, with Llama-3.2-1B as validation.
While the failure modes appear general (confirmed by cross-family consistency), broader model coverage would strengthen the conclusions.

\textbf{SFT $\neq$ pre-training contamination.}
Our contamination injection via LoRA SFT is a proxy for pre-training contamination.
Pre-training contamination may produce different effect signatures (e.g., less catastrophic forgetting due to gradual exposure), potentially making calibration easier---though our identifiability results suggest $N$ and signal quality remain binding constraints regardless.

\textbf{Closed-weight inapplicability.}
ContamCal requires model weights for its most effective signal (S1).
The closed-weight 2-signal mode (S3+S4) is worse than no adjustment, limiting applicability to open-weight models.

\textbf{Frontier scale gap.}
At 72B, all signals are at chance level.
Our conclusions about calibration impossibility at frontier scale rest on extrapolation from smaller models, not direct controlled experiments at 72B.

\textbf{Unidimensional $\theta$.}
We model ability as a single latent dimension, following standard IRT practice.
MMLU spans diverse domains better modeled by multidimensional IRT \cite{reckase2009multidimensional,burnell2023rethink}, which may partially explain MMLU's convergence failures.

%--------------------------------------------------------------
\section{Conclusion}
\label{sec:conclusion}
%--------------------------------------------------------------

We set out to build a principled framework for post-hoc contamination calibration using IRT.
Instead, we discovered why such calibration is structurally infeasible under current evaluation conditions.

Our diagnostic study yields three actionable conclusions for the community:
(1)~Contamination effects are benchmark-dependent, scale-dependent, and non-monotonic---the ``contamination inflates scores'' narrative is an oversimplification that holds only for specific benchmark--scale combinations.
(2)~Detection signals, particularly Min-K++, can actively mislead rather than merely fail; practitioners should validate detection tools on controlled contamination before trusting their output.
(3)~IRT-based calibration requires $\geq$50 independent models and signal AUC$\geq$0.80---conditions that LLM evaluation pipelines do not currently meet and are unlikely to meet in the near term.

Until detection signals improve dramatically and evaluation pipelines include substantially more independent model checkpoints, post-hoc contamination calibration via IRT remains a theoretical possibility, not a practical tool.
The community's best current options remain dynamic benchmarks \cite{white2024livebench}, human evaluation \cite{zheng2024judging}, and transparent reporting of contamination risk rather than statistical correction.

%--------------------------------------------------------------
% References
%--------------------------------------------------------------
%\bibliography{references}
%\bibliographystyle{plainnat}

\begin{thebibliography}{40}

\bibitem[Adams et~al.(1997)]{adams1997multidimensional}
Adams, R.~J., Wilson, M., and Wang, W.-C. (1997).
\newblock The multidimensional random coefficients multinomial logit model.
\newblock \emph{Applied Psychological Measurement}, 21(1):1--23.

\bibitem[Angelopoulos and Bates(2023)]{angelopoulos2023conformal}
Angelopoulos, A.~N. and Bates, S. (2023).
\newblock Conformal prediction: A gentle introduction.
\newblock \emph{Foundations and Trends in Machine Learning}, 16(4):494--591.

\bibitem[Burnell et~al.(2023)]{burnell2023rethink}
Burnell, R., et~al. (2023).
\newblock Rethink the evaluation of large language models with psychometrics.
\newblock \emph{arXiv preprint arXiv:2302.04166}.

\bibitem[Carlini et~al.(2022)]{carlini2022membership}
Carlini, N., et~al. (2022).
\newblock Membership inference attacks from first principles.
\newblock In \emph{IEEE S\&P}, pages 1519--1538.

\bibitem[Carpenter et~al.(2017)]{carpenter2017stan}
Carpenter, B., et~al. (2017).
\newblock Stan: A probabilistic programming language.
\newblock \emph{Journal of Statistical Software}, 76(1).

\bibitem[Clark et~al.(2018)]{clark2018think}
Clark, P., et~al. (2018).
\newblock Think you have solved question answering? Try ARC.
\newblock \emph{arXiv preprint arXiv:1803.05457}.

\bibitem[Cobbe et~al.(2021)]{cobbe2021gsm8k}
Cobbe, K., et~al. (2021).
\newblock Training verifiers to solve math word problems.
\newblock \emph{arXiv preprint arXiv:2110.14168}.

\bibitem[Cormack et~al.(2009)]{cormack2009reciprocal}
Cormack, G.~V., Clarke, C.~L., and Buettcher, S. (2009).
\newblock Reciprocal rank fusion outperforms condorcet and individual rank learning methods.
\newblock In \emph{SIGIR}, pages 758--759.

\bibitem[De~Boeck and Wilson(2004)]{deboeck2004explanatory}
De~Boeck, P. and Wilson, M. (2004).
\newblock \emph{Explanatory Item Response Models}.
\newblock Springer.

\bibitem[Dekoninck et~al.(2024)]{dekoninck2024constat}
Dekoninck, J., et~al. (2024).
\newblock ConStat: Performance-based contamination detection in large language models.
\newblock In \emph{NeurIPS}.

\bibitem[Eckerly(2017)]{eckerly2017detecting}
Eckerly, C.~A. (2017).
\newblock Detecting item preknowledge and item compromise.
\newblock \emph{Journal of Educational Measurement}, 54(4):391--410.

\bibitem[Golchin and Surdeanu(2024)]{golchin2024time}
Golchin, S. and Surdeanu, M. (2024).
\newblock Time travel in LLMs: Tracing data contamination in large language models.
\newblock In \emph{ICLR}.

\bibitem[Guo et~al.(2017)]{guo2017calibration}
Guo, C., et~al. (2017).
\newblock On calibration of modern neural networks.
\newblock In \emph{ICML}, pages 1321--1330.

\bibitem[Hendrycks et~al.(2021)]{hendrycks2021measuring}
Hendrycks, D., et~al. (2021).
\newblock Measuring massive multitask language understanding.
\newblock In \emph{ICLR}.

\bibitem[Holland and Wainer(1993)]{holland1993differential}
Holland, P.~W. and Wainer, H., editors (1993).
\newblock \emph{Differential Item Functioning}.
\newblock Lawrence Erlbaum.

\bibitem[Lalor et~al.(2024)]{lalor2024benchmarking}
Lalor, J.~P., et~al. (2024).
\newblock Benchmarking large language models with item response theory.
\newblock In \emph{AAAI}.

\bibitem[Liang et~al.(2023)]{liang2023holistic}
Liang, P., et~al. (2023).
\newblock Holistic evaluation of language models.
\newblock \emph{TMLR}.

\bibitem[Mislevy(1986)]{mislevy1986bayes}
Mislevy, R.~J. (1986).
\newblock Bayes modal estimation in item response models.
\newblock \emph{Psychometrika}, 51(2):177--195.

\bibitem[Oren et~al.(2024)]{oren2024proving}
Oren, Y., et~al. (2024).
\newblock Proving test set contamination in black-box language models.
\newblock In \emph{EMNLP}.

\bibitem[Platt(1999)]{platt1999probabilistic}
Platt, J.~C. (1999).
\newblock Probabilistic outputs for support vector machines.
\newblock In \emph{Advances in Large Margin Classifiers}, pages 61--74.

\bibitem[Polo et~al.(2024)]{polo2024tinybenchmarks}
Polo, F.~M., et~al. (2024).
\newblock TinyBenchmarks: Evaluating LLMs with fewer examples.
\newblock In \emph{ICML}.

\bibitem[Pruthi et~al.(2020)]{pruthi2020tracin}
Pruthi, G., et~al. (2020).
\newblock Estimating training data influence by tracing gradient descent.
\newblock In \emph{NeurIPS}.

\bibitem[Reckase(2009)]{reckase2009multidimensional}
Reckase, M.~D. (2009).
\newblock \emph{Multidimensional Item Response Theory}.
\newblock Springer.

\bibitem[San~Mart\'in et~al.(2009)]{sanmartin2009identifiability}
San~Mart\'in, E., et~al. (2009).
\newblock Identifiability of parametric item response models.
\newblock \emph{Journal of Statistical Planning and Inference}, 139:2652--2665.

\bibitem[Shi et~al.(2024)]{shi2024detecting}
Shi, W., et~al. (2024).
\newblock Detecting pretraining data from large language models.
\newblock In \emph{ICLR}.

\bibitem[Shokri et~al.(2017)]{shokri2017membership}
Shokri, R., et~al. (2017).
\newblock Membership inference attacks against machine learning models.
\newblock In \emph{IEEE S\&P}, pages 3--18.

\bibitem[Sinharay(2017)]{sinharay2017assessing}
Sinharay, S. (2017).
\newblock Detection of item preknowledge using likelihood ratio test and score test.
\newblock \emph{Journal of Educational and Behavioral Statistics}, 42(1):46--68.

\bibitem[Talts et~al.(2018)]{talts2018sbc}
Talts, S., et~al. (2018).
\newblock Validating Bayesian inference algorithms with simulation-based calibration.
\newblock \emph{arXiv preprint arXiv:1804.06788}.

\bibitem[Thissen et~al.(1993)]{thissen1993detection}
Thissen, D., Steinberg, L., and Wainer, H. (1993).
\newblock Detection of differential item functioning using the parameters of item response models.
\newblock In \emph{Differential Item Functioning}, pages 67--113.

\bibitem[White et~al.(2024)]{white2024livebench}
White, C., et~al. (2024).
\newblock LiveBench: A challenging, contamination-free LLM benchmark.
\newblock \emph{arXiv preprint arXiv:2406.19314}.

\bibitem[Xu et~al.(2024)]{xu2024benchmark}
Xu, R., et~al. (2024).
\newblock Benchmark self-evolving: A multi-agent framework for dynamic LLM evaluation.
\newblock In \emph{ICML}.

\bibitem[Yang et~al.(2023)]{yang2023rethinking}
Yang, S., et~al. (2023).
\newblock Rethinking benchmark and contamination for language models with rephrased samples.
\newblock \emph{arXiv preprint arXiv:2311.04850}.

\bibitem[Zheng et~al.(2024)]{zheng2024judging}
Zheng, L., et~al. (2024).
\newblock Judging LLM-as-a-judge with MT-Bench and Chatbot Arena.
\newblock In \emph{NeurIPS}.

\bibitem[Zumbo(1999)]{zumbo1999handbook}
Zumbo, B.~D. (1999).
\newblock \emph{A Handbook on the Theory and Methods of Differential Item Functioning}.
\newblock NSERC.

\end{thebibliography}

\end{document}
